\documentclass[11pt]{article}
\usepackage{amsmath, amssymb, amsthm}
\usepackage[margin=1in]{geometry}
\usepackage{hyperref}
\usepackage{booktabs}

\newcommand{\R}{\mathbb{R}}
\newcommand{\pd}[2]{\frac{\partial #1}{\partial #2}}
\newcommand{\eps}{\varepsilon}

\title{From the Dirichlet--Neumann Operator to the Nonlinear Schr\"odinger Equation:\\
Deep-Water Wave Regimes and the Benjamin--Feir Instability}
\author{Notes for DNO--FNO project}
\date{April 2026}

\begin{document}
\maketitle

\tableofcontents
\newpage

%=============================================================================
\section{The Water Wave Problem via the DNO}
%=============================================================================

Consider irrotational, incompressible potential flow in the domain
\[
\Omega(t) = \{(x,y) \in \R \times \R : -h < y < \eta(x,t)\},
\]
with free surface $\eta(x,t)$ and flat bottom at $y = -h$.
The velocity potential $\phi(x,y,t)$ satisfies
\[
\Delta \phi = 0 \quad \text{in } \Omega(t), \qquad
\pd{\phi}{y}\bigg|_{y=-h} = 0.
\]
This is a free boundary problem: the domain depends on the solution.

\subsection{Zakharov--Craig--Sulem Formulation}

The key reduction (Zakharov~\cite{zakharov1968}, Craig \& Sulem~\cite{craig1993})
reformulates the problem entirely on the free surface.
Define the surface potential $\xi(x,t) = \phi(x,\eta(x,t),t)$ and the
\textbf{Dirichlet--Neumann operator}:
\[
G(\eta)\xi = \phi_y\big|_{y=\eta} - \eta_x\,\phi_x\big|_{y=\eta}.
\]
This is the standard Dirichlet-to-Neumann map for the Laplacian on $\Omega(t)$,
weighted by the surface geometry.
The Euler equations then reduce to:
\begin{align}
\eta_t &= G(\eta)\xi, \label{eq:kinematic}\\
\xi_t &= -g\eta - \tfrac{1}{2}|\xi_x|^2
+ \frac{\bigl(G(\eta)\xi + \eta_x\,\xi_x\bigr)^2}{2(1+|\eta_x|^2)}.
\label{eq:dynamic}
\end{align}
This is a Hamiltonian system with
$H = \frac{1}{2}\int[\xi\,G(\eta)\xi + g\eta^2]\,dx$
and canonically conjugate variables $(\eta, \xi)$.

The payoff is dimensional reduction: instead of solving a 2D elliptic
problem at every time step, one evaluates $G(\eta)$ acting on a function
of one spatial variable.

\subsection{Taylor Expansion of the DNO}

$G(\eta)$ is analytic in $\eta$ for sufficiently smooth surfaces
(Coifman--Meyer, Craig--Nicholls), with convergent Taylor expansion:
\[
G(\eta) = \sum_{j=0}^{\infty} G_j(\eta),
\qquad G_j \text{ homogeneous of degree } j \text{ in } \eta.
\]
The zeroth-order term is the Fourier multiplier
\begin{equation}
\widehat{G_0\xi}(k) = |k|\tanh(h|k|)\,\hat\xi(k).
\label{eq:G0}
\end{equation}
Higher terms follow from a recursion:
\begin{align}
G_1(\eta) &= D\eta D - G_0\,\eta\,G_0, \\
G_2(\eta) &= -\tfrac{1}{2}\bigl(G_0\,\eta^2 D^2 + D^2\eta^2 G_0
- 2G_0\,\eta\,G_0\,\eta\,G_0\bigr),
\end{align}
where $D = -i\partial_x$.

\paragraph{Depth dependence.}
The symbol $|k|\tanh(h|k|)$ interpolates between two regimes:
\begin{itemize}
\item Finite depth ($h \sim 1$): $G_0$ damps high modes.
\item Deep water ($h \to \infty$): $G_0 \to |D|$, the half-derivative.
\end{itemize}
Since $G_0$ appears in the recursion for every $G_j$, changing $h$ changes
the entire operator.

\paragraph{Numerical truncation.}
In practice, $G(\eta) \approx \sum_{j=0}^{M} G_j(\eta)$ with $M = 4$--$6$.
Each $G_j$ is computed via FFTs at cost $O(M^2 N\log N)$.
Convergence is exponential in $M$ for smooth $\eta$~\cite{xu2009}.


%=============================================================================
\section{Derivation of the NLS}
\label{sec:nls}
%=============================================================================

The \textbf{nonlinear Schr\"odinger equation} (NLS) is the dispersive PDE
\begin{equation}
iA_T + p\,A_{XX} + q\,|A|^2 A = 0,
\label{eq:nls_preview}
\end{equation}
where $A(X,T) \in \mathbb{C}$ is the envelope of a modulated wave train
and $p, q \in \R$ are coefficients determined by the underlying physics.
When $pq > 0$ the equation is \emph{focusing} (admits solitons, exhibits
modulational instability); when $pq < 0$ it is \emph{defocusing}.

In this section we derive \eqref{eq:nls_preview} from the water wave
system \eqref{eq:kinematic}--\eqref{eq:dynamic} in the deep-water limit
$G_0 = |D|$, identifying $p$ and $q$ explicitly in terms of the carrier
wavenumber $k_0$ and the DNO.
The small parameter is the wave steepness $\eps = k_0 a_0$.
We write $E = e^{i\theta}$ throughout.

\subsection{Motivation}

The system \eqref{eq:kinematic}--\eqref{eq:dynamic} is exact but
expensive to simulate and difficult to analyze directly.
In deep water, the dominant observed structures are \emph{modulated
wave trains} --- a fast carrier oscillation whose amplitude and phase
vary slowly in space and time.
The NLS captures the leading-order dynamics of this slow modulation,
reducing the full infinite-dimensional system to \eqref{eq:nls_preview}.

This has both analytical and practical consequences.
Analytically, the NLS is integrable (solvable by inverse scattering),
so its solution space is completely understood: solitons, breathers,
modulational instabilities all have exact descriptions.
Practically, the NLS tells us \emph{which families of surface states
$(\eta, \xi)$ are physically relevant} in deep water --- and hence
what training data a neural operator should see (Section~\ref{sec:datagen}).

\subsection{The Linear Problem}

Linearize \eqref{eq:kinematic}--\eqref{eq:dynamic} around rest
($\eta = \xi = 0$, $G(\eta) \approx G_0$):
$\eta_t = G_0\xi$, $\xi_t = -g\eta$.
Plane waves $\propto e^{i(kx-\omega t)}$ satisfy $\omega^2 = g|k|$.

A narrow-band wave packet centered at $k_0 > 0$ has Fourier support
near $k_0$:
\[
\eta(x,t) = \int_{\R} \hat f(k-k_0)\,e^{i(kx-\omega(k)t)}\,dk
= \biggl(\int \hat f(\eps\kappa)\,e^{i\eps\kappa(x - \omega'(k_0)t)
+ i\frac{\eps^2}{2}\omega''(k_0)\kappa^2 t + \cdots}\,\frac{d\kappa}{\eps}\biggr)\,
e^{i(k_0 x - \omega_0 t)}.
\]
The integral in parentheses is the slowly varying complex envelope $A(X,T)$,
with $X = \eps(x - c_g t)$ and $T = \eps^2 t$.
Its modulus $|A|$ is the local wave amplitude;
$\arg(A)$ is the local phase correction to the carrier $E = e^{i\theta}$.

\subsection{Multiple-Scales Setup}

\paragraph{Fast and slow variables.}
Carrier phase $\theta = k_0 x - \omega_0 t$ with $\omega_0 = \sqrt{gk_0}$.
Slow variables $X = \eps(x - c_g t)$, $T = \eps^2 t$, with group velocity
$c_g = \omega'(k_0) = \omega_0/(2k_0)$.
Derivatives become:
\begin{equation}
\partial_x = k_0\partial_\theta + \eps\partial_X, \qquad
\partial_t = -\omega_0\partial_\theta - \eps c_g\partial_X + \eps^2\partial_T.
\label{eq:deriv_rules}
\end{equation}

\paragraph{$G_0$ on modulated harmonics.}
For $k_0 > 0$, the symbol $|k| = k$ for $k$ near $k_0$ (exactly, not approximately).
So for any slowly modulated harmonic $f(X)\,E^n$ with $n \geq 1$:
\begin{equation}
G_0\bigl[f(X)\,E^n\bigr]
= (nk_0 - i\eps\partial_X)\,f\cdot E^n.
\label{eq:G0_modulated}
\end{equation}
This is \emph{exact} in deep water --- there are no $O(\eps^2)$ corrections
from $G_0$ alone, because $|k|$ is linear for $k > 0$.
(At finite depth, $|k|\tanh(h|k|)$ is not linear, and its curvature
would contribute additional terms here.)

\paragraph{Ansatz.}
Expand $\eta = \eps\eta_1 + \eps^2\eta_2 + \eps^3\eta_3 + \cdots$ and
$\xi = \eps\xi_1 + \eps^2\xi_2 + \eps^3\xi_3 + \cdots$, with
\begin{equation}
\eta_1 = AE + \bar A E^{-1}, \qquad
\xi_1 = BE + \bar B E^{-1},
\label{eq:leading_order}
\end{equation}
where $B(X,T)$ is to be determined in terms of $A$.
The corrections $\eta_j, \xi_j$ for $j \geq 2$ are expanded in Fourier
harmonics: $\eta_j = \sum_n \eta_j^{(n)}(X,T)\,E^n$.


\subsection{$O(\eps)$: Dispersion Relation}

Apply \eqref{eq:deriv_rules} to $\eta_t = G(\eta)\xi$ at $O(\eps)$, coefficient of $E$:
\[
-i\omega_0 A = k_0 B
\quad\implies\quad
B = \frac{-i\omega_0}{k_0}\,A.
\]
From $\xi_t = -g\eta$ at $O(\eps)$, coefficient of $E$:
\[
-i\omega_0 B = -gA
\quad\implies\quad
B = \frac{ig}{\omega_0}\,A = \frac{-i\omega_0}{k_0}\,A,
\]
where the last equality uses $g = \omega_0^2/k_0$.
Both equations give the same $B$; consistency requires
$\omega_0^2 = gk_0$. \hfill$\square$


\subsection{$O(\eps^2)$: Group Velocity and Forced Harmonics}

At this order, the slow $X$-derivatives of $A, B$ enter, and the
quadratic nonlinearities in \eqref{eq:dynamic} generate forcing at
harmonics $E^0$ (mean) and $E^2$ (second harmonic).

\subsubsection{No secularity at $E^1$}

The system at the fundamental $E^1$ reads (kinematic, then dynamic):
\begin{align}
-i\omega_0\eta_2^{(1)} - c_g A_X
&= k_0\xi_2^{(1)} + (-iB_X), \label{eq:K2_E1}\\[4pt]
-i\omega_0\xi_2^{(1)} - c_g B_X
&= -g\eta_2^{(1)}. \label{eq:D2_E1}
\end{align}
The $-iB_X$ in \eqref{eq:K2_E1} comes from
the $\eps\partial_X$ correction in \eqref{eq:G0_modulated};
the $-c_g A_X$ and $-c_g B_X$ come from
$\eps c_g\partial_X$ in $\partial_t$.
(No nonlinear terms at $E^1$:
the quadratic products $\xi_{1,x}^2$ and $(G_0\xi_1)^2$ produce
only $E^0$ and $E^2$.)

Equations \eqref{eq:K2_E1}--\eqref{eq:D2_E1} are two equations for
$\eta_2^{(1)}$ and $\xi_2^{(1)}$, but the coefficient matrix
$\bigl(\begin{smallmatrix} -i\omega_0 & -k_0 \\
-g & i\omega_0 \end{smallmatrix}\bigr)$
is singular ($\omega_0^2 = gk_0$).
A solution exists iff the forcing satisfies a solvability condition.
Substituting $B = -i\omega_0 A/k_0$ and $c_g = \omega_0/(2k_0)$
into the forcing terms, one checks directly that the $A_X$ terms cancel.
So there is no secularity at $E^1$ --- the group velocity choice
$c_g = \omega'(k_0)$ guarantees this.
We set $\eta_2^{(1)} = \xi_2^{(1)} = 0$ (absorbed into $A$ by redefinition).

\subsubsection{Second harmonic $E^2$}

The quadratic nonlinearities contribute at $E^2$.
Compute the relevant products at leading order, using $B = -i\omega_0 A/k_0$:
\begin{align}
\xi_{1,x} &= ik_0 BE - ik_0\bar B E^{-1}, \\
\xi_{1,x}^2 &= -k_0^2\bigl(B^2 E^2 - 2|B|^2 + \bar B^2 E^{-2}\bigr), \\
(G_0\xi_1)^2 &= k_0^2\bigl(B^2 E^2 + 2|B|^2 + \bar B^2 E^{-2}\bigr).
\end{align}
The combined nonlinear forcing in \eqref{eq:dynamic} at $E^2$ is
\[
\bigl[-\tfrac{1}{2}\xi_{1,x}^2 + \tfrac{1}{2}(G_0\xi_1)^2\bigr]_{E^2}
= k_0^2 B^2 = -\omega_0^2 A^2.
\]
(The numerator $(G\xi + \eta_x\xi_x)^2$ reduces to $(G_0\xi_1)^2$
at this order because $\eta_x\xi_x$ is $O(\eps^2)$ but only has
$E^0$ and $E^2$ --- it contributes the same way.)

The kinematic and dynamic equations at $E^2$ give:
\begin{align}
-2i\omega_0\eta_2^{(2)} &= 2k_0\xi_2^{(2)}, \label{eq:K2_E2}\\
-2i\omega_0\xi_2^{(2)} &= -g\eta_2^{(2)} - \omega_0^2 A^2. \label{eq:D2_E2}
\end{align}
From \eqref{eq:K2_E2}: $\xi_2^{(2)} = -i\omega_0\eta_2^{(2)}/k_0$.
Substitute into \eqref{eq:D2_E2}:
\[
-2i\omega_0\cdot\frac{-i\omega_0\eta_2^{(2)}}{k_0} = -g\eta_2^{(2)} - \omega_0^2 A^2
\quad\implies\quad
-\frac{2\omega_0^2}{k_0}\eta_2^{(2)} = -g\eta_2^{(2)} - \omega_0^2 A^2.
\]
Using $g = \omega_0^2/k_0$:
the left side is $-2g\eta_2^{(2)}$, the right is $-g\eta_2^{(2)} - \omega_0^2 A^2$,
giving
\begin{equation}
\boxed{\eta_2^{(2)} = k_0 A^2}, \qquad
\xi_2^{(2)} = \frac{-i\omega_0}{k_0}\cdot k_0 A^2 = -i\omega_0 A^2.
\label{eq:eta2_E2}
\end{equation}

\subsubsection{Mean flow $E^0$}

The $E^0$ nonlinear forcing in the dynamic equation is
\[
\bigl[-\tfrac{1}{2}\xi_{1,x}^2 + \tfrac{1}{2}(G_0\xi_1)^2\bigr]_{E^0}
= k_0^2|B|^2 - k_0^2|B|^2 = 0.
\]
This cancellation is specific to deep water (at finite depth, $G_0$
on $E^{+1}$ and $E^{-1}$ gives different symbols, and the mean flow
is nonzero).  Therefore:
\begin{equation}
\eta_2^{(0)} = 0, \qquad \xi_2^{(0)} = 0 \quad\text{(deep water)}.
\label{eq:eta2_E0}
\end{equation}

\subsection{$O(\eps^3)$: The NLS}

\subsubsection{Solvability framework}

At $O(\eps^3)$, the system at $E^1$ has the same singular matrix as before:
\begin{align}
-i\omega_0\eta_3^{(1)} - k_0\xi_3^{(1)} &= R_K, \label{eq:sys3_K}\\
-g\eta_3^{(1)} + i\omega_0\xi_3^{(1)} &= R_D, \label{eq:sys3_D}
\end{align}
where $R_K, R_D$ collect all $O(\eps^3)$, $E^1$ terms not involving
$\eta_3^{(1)}$ or $\xi_3^{(1)}$.
The solvability condition (left null vector of the coefficient matrix) is:
\begin{equation}
i\omega_0\,R_K + k_0\,R_D = 0.
\label{eq:solvability}
\end{equation}
We now compute $R_K$ and $R_D$.

\subsubsection{Dispersive contribution (the $p\,A_{XX}$ term)}

The dispersion comes from the interplay of $\partial_t^2$ and $G_0$
across the kinematic and dynamic equations, not from $G_0$ alone.
The cleanest way to see this: differentiate $\eta_t = G_0\xi + \cdots$
in time to get $\eta_{tt} = -gG_0\eta + \cdots$, i.e.\
$(\partial_t^2 + gG_0)\eta = \text{nonlinear forcing}$.

Acting on $\eps A(X,T)\,E$:
\begin{align}
\partial_t(\eps AE) &= \eps(-i\omega_0 A - \eps c_g A_X + \eps^2 A_T)\,E,
\\
\partial_t^2(\eps AE) &= \eps\bigl(-\omega_0^2 A
+ 2i\omega_0\eps c_g A_X
+ \eps^2(c_g^2 A_{XX} - 2i\omega_0 A_T)\bigr)\,E,
\\
gG_0(\eps AE) &= \eps\,(gk_0 A - ig\eps A_X)\,E
= \eps\,(\omega_0^2 A - i\eps\omega_0^2 A_X/k_0)\,E.
\end{align}
Adding:
\[
(\partial_t^2 + gG_0)[\eps AE]
= \eps^3\bigl(c_g^2 A_{XX} - 2i\omega_0 A_T\bigr)\,E
\]
--- the $O(\eps)$ and $O(\eps^2)$ terms cancel
(using $\omega_0^2 = gk_0$ and $c_g = \omega_0/(2k_0)$).

So the linear part of the solvability condition at $O(\eps^3)$ contributes
$c_g^2 A_{XX} - 2i\omega_0 A_T$.
Rearranging:
\[
iA_T = \frac{c_g^2}{2\omega_0}\,A_{XX} + \cdots
= \frac{\omega_0}{8k_0^2}\,A_{XX} + \cdots
\]
This gives the dispersive coefficient:
\begin{equation}
p = \frac{c_g^2}{2\omega_0} = \frac{\omega_0}{8k_0^2}.
\label{eq:p}
\end{equation}

\subsubsection{Nonlinear contribution (the $q\,|A|^2A$ term)}

The nonlinear forcing at $O(\eps^3)$, $E^1$ has several sources.
We enumerate them, noting that $\eta_2^{(0)} = 0$ (deep water) eliminates
several terms.

\paragraph{(a) Quadratic terms in \eqref{eq:dynamic} coupling $\eta_1$ with $\eta_2$.}
The products $\xi_{1,x}\cdot\xi_{2,x}$ and $G_0\xi_1\cdot G_0\xi_2$
at $E^1$ involve
$\xi_1$ ($E^{\pm 1}$) interacting with $\xi_2^{(2)} E^2$:
\begin{align}
[\xi_{1,x}\cdot\xi_{2,x}]_{E^1}
&= (ik_0\bar B)(-2ik_0)(-i\omega_0 A^2)
+ \text{(from } E^{-1}\cdot E^2\text{)}, \\
[G_0\xi_1\cdot G_0\xi_2]_{E^1}
&= (k_0\bar B)(2k_0)(-i\omega_0 A^2).
\end{align}

\paragraph{(b) $G_1(\eta_2)\xi_1$ at $E^1$.}
$G_1(\eta)\xi = D(\eta\,D\xi) - G_0(\eta\,G_0\xi)$.
With $\eta_2^{(2)}E^2$ and $\xi_1 = BE + \bar BE^{-1}$:
the $E^1$ component comes from $E^2 \cdot E^{-1}$, giving terms
proportional to $\eta_2^{(2)}\bar B = k_0 A^2 \cdot (i\omega_0\bar A/k_0)$.

\paragraph{(c) $G_1(\eta_1)\xi_2$ at $E^1$.}
Similarly, $\eta_1 = AE + \bar AE^{-1}$ with
$\xi_2^{(2)}E^2$: the $E^1$ component comes from $E^{-1}\cdot E^2$,
proportional to $\bar A \cdot \xi_2^{(2)} = \bar A(-i\omega_0 A^2)$.

\paragraph{(d) $G_2(\eta_1)\xi_1$ at $E^1$.}
This involves $\eta_1^2$ (which has $E^0$ and $E^{\pm 2}$ components)
acting on $\xi_1$ through the second-order DNO correction.
The $E^1$ component comes from the $E^0$ part of $\eta_1^2$ ($= 2|A|^2$)
times the $E^1$ part of $\xi_1$.

\medskip

Each of (a)--(d) produces a term proportional to $|A|^2 A\cdot E$.
Collecting all contributions and applying the solvability condition
\eqref{eq:solvability}, the nonlinear coefficient works out to:
\begin{equation}
q = \frac{\omega_0 k_0^2}{2}.
\label{eq:q}
\end{equation}
(The detailed algebra is carried out in, e.g.,
Mei, Stiassnie \& Yue~\cite{mei2005}, \S13.3,
or Ablowitz~\cite{ablowitz2011}, \S6.4.)

\subsubsection{The NLS}

Combining the dispersive and nonlinear contributions, the
solvability condition at $O(\eps^3)$ is
\begin{equation}
\boxed{
iA_T + p\,A_{XX} + q\,|A|^2 A = 0
}
\label{eq:nls}
\end{equation}
with
\[
p = \frac{\omega_0}{8k_0^2}, \qquad
q = \frac{\omega_0 k_0^2}{2}.
\]
Since $pq > 0$ in deep water, the equation is \textbf{focusing}:
it admits soliton solutions, and plane waves are modulationally unstable.

\paragraph{Rigorous justification.}
D\"ull, Schneider \& Wayne~\cite{dull2015} proved that NLS solutions
approximate the full Euler equations on the natural $O(1/\eps^2)$ time scale,
with the key estimate being analyticity bounds on $G(\eta)$.

\subsection{Finite Depth}

For finite $h$, the same procedure applies but $G_0 = |D|\tanh(h|D|)$,
so \eqref{eq:G0_modulated} acquires $O(\eps^2)$ corrections from the
curvature of the symbol.
The dispersion relation becomes $\omega = \sqrt{gk\tanh(hk)}$,
and critically, $\eta_2^{(0)} \neq 0$ --- the mean flow no longer vanishes,
modifying the nonlinear coefficient.
The sign of $q$ changes at $k_0 h \approx 1.363$: below this threshold,
$pq < 0$ (defocusing), and plane waves are stable.


%=============================================================================
\section{What the NLS--Euler Connection Buys Us}
\label{sec:connection}
%=============================================================================

The derivation in Section~\ref{sec:nls} establishes that the NLS
\eqref{eq:nls} governs the envelope dynamics of deep-water waves
at $O(\eps^3)$.
This is useful in two distinct ways: the NLS gives us a
\emph{catalog of solutions}, and the Euler equations (via the DNO)
give us the \emph{exact nonlinear operator} we want to learn.

\subsection{The NLS as a Solution Generator}

The focusing NLS is integrable: its full solution space is known
via inverse scattering.
This gives closed-form expressions for envelope solitons, Akhmediev
breathers, Peregrine breathers, and multi-soliton states ---
none of which are easy to guess from the Euler equations directly.

Each NLS solution $A(X,T)$ lifts to physical variables $(\eta,\xi)$
via \eqref{eq:leading_order}, producing an approximate solution of the
full water wave equations valid to $O(\eps^2)$.
These lifted states serve as:
\begin{itemize}
\item \textbf{Initial conditions} for DNO-based time integration of the
full Euler equations.
\item \textbf{Standalone data points}: evaluate $G(\eta)\xi$ directly
on the lifted $(\eta,\xi)$ pair using the DNO series, without any
time integration at all.
\end{itemize}

\subsection{What Is Easy and What Is Not}

\paragraph{Envelope solitons: easy.}
The NLS soliton \eqref{eq:envelope_soliton} gives an explicit formula
for $A$ at any $(X,T)$.
Lifting to $(\eta,\xi)$ is a single evaluation of
\eqref{eq:leading_order} --- no time stepping, no iteration.
One can generate thousands of soliton snapshots at different $(k_0, A_0)$
and evaluate $G(\eta)\xi$ on each.
The only cost is the DNO series evaluation itself.

\paragraph{Peregrine and Akhmediev breathers: easy as snapshots, interesting as trajectories.}
Like the soliton, the Peregrine \eqref{eq:peregrine} and Akhmediev
\eqref{eq:akhmediev} breathers have closed-form expressions.
Generating $(\eta,\xi)$ snapshots at any $(X,T)$ is immediate.
However, these are approximate solutions of Euler --- the NLS truncates
at $O(\eps^3)$, so the lifted state drifts from the true Euler
trajectory over time.
If one wants data that reflects the \emph{full} nonlinear dynamics
(post-NLS corrections, recurrence breakdown, frequency downshift),
one must time-integrate the Euler equations from the NLS-derived
initial condition and snapshot along the trajectory.

\paragraph{Benjamin--Feir recurrence: requires time integration.}
The BF instability is easy to \emph{set up} --- perturb a Stokes
wave with sidebands, which the NLS predicts will grow.
But the interesting physics (recurrence, eventual breakdown) plays out
over hundreds of carrier periods and involves the full nonlinear DNO.
The NLS predicts exact recurrence; the Euler equations show it
breaking down after 2--3 cycles~\cite{xu2009}.
Capturing this discrepancy requires running the DNO solver, which is
expensive but produces the most physically rich training data.

\subsection{Summary: Division of Labor}

\begin{center}
\begin{tabular}{@{}lll@{}}
\toprule
\textbf{Role} & \textbf{NLS provides} & \textbf{Euler/DNO provides} \\
\midrule
Solution families & closed-form catalog & --- \\
Initial conditions & $(\eta,\xi)$ via lifting & --- \\
$G(\eta)\xi$ evaluation & --- & exact (to order $M$) \\
Long-time dynamics & approximate (integrable) & exact (non-integrable) \\
Post-NLS effects & --- & recurrence breakdown, \\
 & & frequency downshift, breaking \\
\bottomrule
\end{tabular}
\end{center}

In short: the NLS tells us \emph{what to simulate}, and the DNO tells
us \emph{what the operator actually does} on those states.
For operator learning, the NLS is a recipe book for generating
diverse, physically motivated inputs $(\eta,\xi)$; the DNO series
provides the ground-truth labels $G(\eta)\xi$.


%=============================================================================
\section{Exact Solutions of the Focusing NLS}
\label{sec:solutions}
%=============================================================================

The focusing NLS is integrable (inverse scattering), yielding
several families of exact solutions that correspond to distinct
deep-water wave phenomena.
Each can be lifted to $(\eta, \xi)$ via
\eqref{eq:leading_order}.

\subsection{Plane Wave (Stokes Wave)}
\begin{equation}
A(X,T) = A_0\,e^{iqA_0^2 T}.
\end{equation}
Spatially uniform carrier with nonlinear frequency correction.
Starting point for stability analysis.

\subsection{Envelope Soliton}
\begin{equation}
A(X,T) = A_0\,\operatorname{sech}\!\Bigl(A_0\sqrt{\tfrac{q}{2p}}\;X\Bigr)
\exp\!\Bigl(\tfrac{i}{2}qA_0^2 T\Bigr).
\label{eq:envelope_soliton}
\end{equation}
A localized wave packet propagating at $c_g$ without dispersing ---
the deep-water counterpart of the KdV soliton, but modulating a carrier
rather than forming a single hump.
At $T = 0$:
\begin{equation}
\eta(x,0) = 2\eps A_0\,
\operatorname{sech}\!\bigl(\eps A_0\sqrt{q/(2p)}\;x\bigr)\cos(k_0 x).
\label{eq:eta_soliton}
\end{equation}
Confirmed experimentally~\cite{shemer2013}.

\subsection{Akhmediev Breather}

Spatially periodic, temporally localized~\cite{akhmediev1987}:
\begin{equation}
A(X,T) = A_0\,\frac{\cosh(\sigma T - 2i\varphi)
- \cos\varphi\cos(\kappa X)}
{\cosh(\sigma T) - \cos\varphi\cos(\kappa X)}\,e^{iA_0^2 T},
\label{eq:akhmediev}
\end{equation}
where $\cos\varphi = \kappa/(2A_0\sqrt{q/p})$ and
$\sigma = \kappa\sqrt{4qA_0^2/p - \kappa^2 q/p}$.

As $T \to -\infty$ this is the plane wave plus an infinitesimal
$\cos(\kappa X)$ perturbation.
At $T = 0$ the modulation reaches maximum amplitude.
As $T \to +\infty$ it returns to the plane wave.
This describes one complete cycle of the Benjamin--Feir instability
(Section~\ref{sec:bf}).

\subsection{Peregrine Breather}

The $\kappa \to 0$ limit of the Akhmediev breather~\cite{peregrine1983}:
\begin{equation}
A(X,T) = A_0\left(1 - \frac{4(1 + 2iqA_0^2 T)}
{1 + 4q^2 A_0^4 T^2/p^2 + 2qA_0^2 X^2/p}\right)e^{iqA_0^2 T}.
\label{eq:peregrine}
\end{equation}
Localized in both space and time, with maximum amplitude $3A_0$ ---
the prototypical rogue wave.
Produces steep, localized surface states where $G_2, G_3, \ldots$
are non-negligible.


%=============================================================================
\section{The Benjamin--Feir Instability}
\label{sec:bf}
%=============================================================================

Benjamin and Feir~\cite{benjamin1967} showed that uniform deep-water
Stokes waves are modulationally unstable.
This is a consequence of the focusing character of the NLS.

\subsection{Linear Stability Analysis}

Perturb the plane wave at sideband wavenumbers $k_0 \pm \delta k$:
\[
A = \bigl(A_0 + a_+\,e^{i\delta k\,X + \Omega T}
+ a_-\,e^{-i\delta k\,X + \bar\Omega T}\bigr)\,e^{iqA_0^2 T}.
\]
Linearizing \eqref{eq:nls}:
\begin{equation}
\Omega^2 = p^2(\delta k)^2\bigl[(\delta k)^2 - 2qA_0^2/p\bigr].
\label{eq:bf_growth}
\end{equation}
Instability ($\Omega^2 < 0$) when
\begin{equation}
0 < |\delta k| < A_0\sqrt{2q/p} = \sqrt{2}\,k_0^2 a_0.
\label{eq:bf_criterion}
\end{equation}
Maximum growth rate $\Omega_{\max} = qA_0^2 = \omega_0 k_0^2 a_0^2/2$
at $|\delta k| = A_0\sqrt{q/p}$.
At steepness $k_0 a_0 = 0.13$, the $e$-folding time is $\sim 19$ carrier periods.

The coefficient $q$ --- and hence the instability bandwidth --- originates
from the nonlinear interaction mediated by $G_1(\eta)\xi$.
The exact stability analysis (Berti, Maspero \& Ventura~\cite{berti2022})
uses the shape derivative of the full DNO.

\subsection{Nonlinear Recurrence}

The instability saturates.
Energy cycles between carrier and sidebands in Fermi--Pasta--Ulam recurrence:
growth $\to$ sideband maximum $\to$ return to carrier $\to$ repeat.
The Akhmediev breather \eqref{eq:akhmediev} is the exact NLS solution
describing one cycle.

For the full Euler equations (via DNO), recurrence persists for 2--3 cycles
before higher-order effects (absent at $O(\eps^3)$) disrupt it.

\subsection{DNO Simulation (Xu \& Guyenne 2009)}

Xu \& Guyenne~\cite{xu2009} simulated this with the full system
\eqref{eq:kinematic}--\eqref{eq:dynamic}, $M = 4$.
Initial condition: Stokes wave ($k_1 = 9$, $k_1 a = 0.13$) with
sideband perturbations at $k_l = 7$, $k_r = 11$:
\begin{align}
\eta(x,0) &= \eta_0 + \eps a\left[\cos\!\left(k_l x - \tfrac{\pi}{4}\right)
+ \cos\!\left(k_r x - \tfrac{\pi}{4}\right)\right],
\label{eq:bf_eta_ic}\\
\xi(x,0) &= \xi_0 + \eps a\left[\frac{e^{k_l\eta}}{\sqrt{k_l}}
\sin\!\left(k_l x - \tfrac{\pi}{4}\right)
+ \frac{e^{k_r\eta}}{\sqrt{k_r}}
\sin\!\left(k_r x - \tfrac{\pi}{4}\right)\right],
\label{eq:bf_xi_ic}
\end{align}
where $(\eta_0, \xi_0)$ is a numerically exact Stokes wave and $\eps = 0.1$.

Results: three full recurrence cycles over 716 carrier periods.
Energy conserved to $O(10^{-3}\%)$ for $t/T \leq 600$.
$M = 2$ suffices (Benjamin--Feir is a four-wave process).

\subsection{Post-NLS Effects}

The full Euler equations exhibit phenomena absent from NLS:
\begin{itemize}
\item Recurrence breakdown after 2--3 cycles (higher-order resonances).
\item Permanent frequency downshift ($O(\eps^4)$ effect).
\item Wave breaking for $k_0 a \gtrsim 0.3$ (the DNO ceases to be defined
when $\eta$ becomes multi-valued).
\end{itemize}
These are precisely the regimes where a learned approximation of the
full DNO would be valuable: it captures the complete nonlinear map
without truncation at finite $M$.


%=============================================================================
\section{Generating Deep-Water Training Data}
\label{sec:datagen}
%=============================================================================

The preceding analysis provides five families of physically consistent
$(\eta, \xi)$ pairs in deep water.
For each, one evaluates $G(\eta)\xi$ via the DNO series at $h = 1000$
(deep-water proxy).

\begin{center}
\begin{tabular}{@{}llp{5cm}@{}}
\toprule
\textbf{Source} & \textbf{Parameters} & \textbf{Nonlinearity} \\
\midrule
Stokes waves & $k_0, a_0$ & mild (single frequency) \\
Linear superposition & random spectrum & nearly linear, broad bandwidth \\
Envelope solitons & $A_0, k_0$ & moderate (localized envelope) \\
BF trajectories & $k_0, a_0, \delta k, t$ & strong (time-varying cascades) \\
Akhmediev/Peregrine & $\kappa, A_0$ & extreme (up to $3\times$ background) \\
\bottomrule
\end{tabular}
\end{center}

\paragraph{Linear superposition} is cheapest:
draw random $k_j, a_j, \phi_j$ and set
\[
\eta(x) = \textstyle\sum_j a_j\cos(k_j x + \phi_j), \qquad
\xi(x) = \textstyle\sum_j \frac{a_j\omega_j}{k_j}\sin(k_j x + \phi_j),
\quad \omega_j = \sqrt{g|k_j|}.
\]

\paragraph{BF trajectories} are most informative:
time-integrate \eqref{eq:kinematic}--\eqref{eq:dynamic} from
\eqref{eq:bf_eta_ic}--\eqref{eq:bf_xi_ic} and snapshot at intervals.

\paragraph{Depth conditioning.}
A model trained at $h = 1$ learns $G_0 = |k|\tanh(k)$ and cannot
extrapolate to $G_0 = |k|$.
For finite depths, the rescaling $x \to x/h$, $\eta \to \eta/h$
reduces to $h = 1$.
The deep-water limit is genuinely different and requires either direct
training data or a depth-conditioned architecture (e.g.\ encoding
$\tanh(hk)$ as an input channel).


%=============================================================================
\begin{thebibliography}{99}

\bibitem{zakharov1968}
V.E.\ Zakharov,
``Stability of periodic waves of finite amplitude on the surface of a deep fluid,''
\textit{J.\ Appl.\ Mech.\ Tech.\ Phys.}\ \textbf{9}(2), 190--194 (1968).

\bibitem{craig1993}
W.\ Craig \& C.\ Sulem,
``Numerical simulation of gravity waves,''
\textit{J.\ Comput.\ Phys.}\ \textbf{108}(1), 73--83 (1993).

\bibitem{craig1997}
W.\ Craig, U.\ Schanz \& C.\ Sulem,
``The modulational regime of three-dimensional water waves and the
Davey--Stewartson system,''
\textit{Ann.\ Inst.\ H.\ Poincar\'e (C)}\ \textbf{14}(5), 615--667 (1997).

\bibitem{craig2004}
W.\ Craig \& C.\ Sulem,
``Asymptotic nonlinear wave modeling through the Dirichlet-to-Neumann operator,''
\textit{Methods Appl.\ Anal.}\ \textbf{11}(4), 475--492 (2004).

\bibitem{dull2015}
W.-P.\ D\"ull, G.\ Schneider \& C.E.\ Wayne,
``Justification of the nonlinear Schr\"odinger equation for the evolution
of gravity driven 2D surface water waves,''
\textit{Arch.\ Ration.\ Mech.\ Anal.}\ \textbf{217}(3), 1219--1270 (2015).

\bibitem{guyenne2007}
P.\ Guyenne \& D.P.\ Nicholls,
``Comparison of Dirichlet--Neumann operator expansions for nonlinear
surface gravity waves,''
\textit{Coastal Eng.}\ \textbf{55}(2), 120--133 (2008).

\bibitem{xu2009}
L.\ Xu \& P.\ Guyenne,
``Numerical simulation of three-dimensional nonlinear water waves,''
\textit{J.\ Comput.\ Phys.}\ \textbf{228}(22), 8446--8466 (2009).

\bibitem{berti2022}
M.\ Berti, A.\ Maspero \& P.\ Ventura,
``Full description of Benjamin--Feir instability of Stokes waves in deep water,''
\textit{Invent.\ Math.}\ \textbf{230}, 651--711 (2022).

\bibitem{benjamin1967}
T.B.\ Benjamin \& J.E.\ Feir,
``The disintegration of wave trains on deep water. Part 1. Theory,''
\textit{J.\ Fluid Mech.}\ \textbf{27}(3), 417--430 (1967).

\bibitem{peregrine1983}
D.H.\ Peregrine,
``Water waves, nonlinear Schr\"odinger equations and their solutions,''
\textit{J.\ Austral.\ Math.\ Soc.\ Ser.\ B}\ \textbf{25}(1), 16--43 (1983).

\bibitem{shemer2013}
L.\ Shemer, A.\ Sergeeva \& A.\ Slunyaev,
``Simulations and experiments of short intense envelope solitons of
surface water waves,''
\textit{Phys.\ Fluids}\ \textbf{25}, 067105 (2013).

\bibitem{akhmediev1987}
N.N.\ Akhmediev, V.M.\ Eleonskii \& N.E.\ Kulagin,
``Exact first-order solutions of the nonlinear Schr\"odinger equation,''
\textit{Theor.\ Math.\ Phys.}\ \textbf{72}, 809--818 (1987).

\bibitem{mei2005}
C.C.\ Mei, M.\ Stiassnie \& D.K.-P.\ Yue,
\textit{Theory and Applications of Ocean Surface Waves},
World Scientific (2005).

\bibitem{ablowitz2011}
M.J.\ Ablowitz,
\textit{Nonlinear Dispersive Waves: Asymptotic Analysis and Solitons},
Cambridge University Press (2011).

\end{thebibliography}

\end{document}
